\subsection{Criterio de decisión y versión 1}

Con $\rho = 0{,}7613$ las prioridades muestran una tendencia común. El núcleo Must (R01, R04, R06, R07, R10, R14, R17 y R19) se decide con la tabla de ponderación entregada, el criterio clínico y la criticidad del dominio. Kano cubre solo ocho características y no midió la disponibilidad (R14) ni la recuperación (R19); por eso no se usa para respaldar esos dos requisitos. El triaje aparece como K3, sin requisito en el catálogo votado, y se documenta como brecha en lugar de incorporarlo sin votación. Se decide:

\begin{enumerate}
  \item \textbf{Incluir en la versión 1} los ocho Must: R01, R04, R06, R07, R10, R14, R17 y R19. Su esfuerzo es 114 h (57,0\%), dentro del límite del 60\%.
  \item \textbf{Incluir condicionalmente} los Should R02, R09, R12, R13 y R15 (34 h) si la capacidad restante lo permite; R05 (12 h) se adopta como restricción tecnológica, no como función condicional. El grupo Should de la tabla entregada suma 46 h. R04 es el registro manual obligatorio del flujo de atención, mientras que R13 es la automatización en quiosco con impresión de tickets y depende de equipo adicional: son funciones distintas, R13 no sustituye a R04 y por eso queda en segundo lugar dentro del grupo.
  \item \textbf{Promover a revisión prioritaria} R03 (18 puntos) y R18 (14 puntos). R03 empata con R06 y R19 y supera a R04 (13); R18 supera solo a R04 dentro de los Must. Ambos están como Could y requieren validación antes de pasar a Should.
  \item \textbf{Respetar las restricciones del proyecto}: R11 (entrega en seis meses) y R20 (revisión de código antes de integrar) no son funcionalidades diferibles; se cumplen aunque R20 figure como Won't have en la tabla de ponderación entregada, donde se registra la contradicción. R05 (pila tecnológica) y R16 (validación quincenal) también son restricciones, no entregables opcionales. Como función opcional queda fuera de la versión 1 solo R08 (compatibilidad Android 10), por su puntaje y su carácter de compatibilidad.
  \item \textbf{Registrar brechas Kano}: K3 (triaje), K7 (calificación de la atención) y K8 (medicamentos recetados) son categorías M/O sin correlato en el catálogo votado. No se incorporan a la versión 1 sin elicitar y votar los requisitos correspondientes. K2 (recordatorios) y K6 (accesibilidad) ya tienen correlato en R02 y R10; K5 (módulos) es indiferente.
\end{enumerate}

Lista final ordenada de mayor a menor total: R07; R01, R10, R14, R17; R03, R06, R19; R02, R12; R18; R04; R09, R13, R15; R05, R08, R11, R16, R20. Los grupos separados por punto y coma comparten rango según la Tabla~\ref{tab:ranking_votacion}.

La decisión se apoya en la Tabla~\ref{tab:requisitos}, la Tabla~\ref{tab:esfuerzo}, los resultados Kano de la Tabla~\ref{tab:resultado_kano}, el orden Kano de la Tabla~\ref{tab:orden_kano}, la votación de la Tabla~\ref{tab:ranking_votacion} y los órdenes comparados de la Tabla~\ref{tab:ordenes}, sin introducir requisitos ni participantes que no estén en los datos entregados. El catálogo tiene 10 funcionales, 6 de calidad y 4 restricciones, por debajo de los 12 y 8 previstos en el kit; la diferencia se declara como limitación y no se cubre con requisitos inventados.

\emph{Transparencia de la votación}: el manifiesto \texttt{evidencias/MANIFIESTO.md} reporta que en 15 de los 20 requisitos el valor de V2 coincide con el promedio aritmético de V3 y V4. Las cifras transcritas no se modificaron para disimular ese patrón; las tres boletas originales se custodian fuera del repositorio para proteger los datos de las personas participantes y se mostrarán al docente si las solicita. Esta limitación se declara sin afirmar una independencia que no puede verificarse desde el repositorio.
